% Paper 1: Senjam et al. (Clinical Ophthalmology, 2021)
\begin{frame}{Literature Review: Smartphones-Based Assistive Technology}
\begin{columns}[T]
  \begin{column}{0.31\textwidth}
    \begin{block}{Paper Overview \cite{senjam2021smartphones}}
      \textbf{Authors:} S. S. Senjam et al.\\
      \textbf{Journal:} Clinical Ophthalmology (2021)\\
      \textbf{Focus:} Evaluation of built-in smartphone features and assistive apps for visually impaired persons.
    \end{block}
    \vspace{0.1cm}
    \begin{block}{Main Contributions}
      \begin{itemize}
        \item Categorised smartphone accessibility apps for vision assistance.
        \item Evaluated screen readers, magnification, and AI scene description.
      \end{itemize}
    \end{block}
  \end{column}
    
  \begin{column}{0.34\textwidth}
    \begin{block}{Methodology \& Technology}
      \begin{itemize}
        \item Smartphone-based computer vision and sensor fusion.
        \item Evaluated apps like Seeing AI, Envision, and Google Lookout.
        \item Audio/speech interface for feedback.
      \end{itemize}
    \end{block}
    \vspace{0.1cm}
    \begin{block}{Inputs \& Outputs}
      \textbf{Input:} RGB camera stream, touch gestures, voice commands.\\
      \textbf{Output:} Synthesized speech, haptic cues.
    \end{block}
  \end{column}
  
  \begin{column}{0.31\textwidth}
    \begin{block}{Summary of Results}
      \begin{itemize}
        \item Over 70\% improvement in independent daily task execution.
        \item Validated smartphone as a primary assistive medium.
      \end{itemize}
    \end{block}
    \vspace{0.1cm}
    \begin{alertblock}{Key Limitations}
      \begin{itemize}
        \item Limited real-time navigation guidance in unstructured road settings.
        \item High dependency on cloud network connectivity.
      \end{itemize}
    \end{alertblock}
  \end{column}
\end{columns}
\end{frame}

% Paper 2: Kathiria et al. (Neurocomputing, 2024)
\begin{frame}{Literature Review: Assistive Systems for Visually Impaired People}
\begin{columns}[T]
  \begin{column}{0.31\textwidth}
    \begin{block}{Paper Overview \cite{kathiria2024assistive}}
      \textbf{Authors:} P. Kathiria et al.\\
      \textbf{Journal:} Neurocomputing (2024)\\
      \textbf{Focus:} Comprehensive survey of advancements and state-of-the-art in assistive navigation systems.
    \end{block}
    \vspace{0.1cm}
    \begin{block}{Main Contributions}
      \begin{itemize}
        \item Benchmark of sensors: LiDAR vs. Ultrasound vs. Cameras.
        \item Taxonomy of navigation, hazard detection, and user interfaces.
      \end{itemize}
    \end{block}
  \end{column}
    
  \begin{column}{0.34\textwidth}
    \begin{block}{Methodology \& Technology}
      \begin{itemize}
        \item Deep learning object detectors (YOLO, SSD, Faster R-CNN).
        \item Multimodal sensory integration (RGB-D + IMU + GPS).
        \item Edge computing vs. Cloud inference trade-offs.
      \end{itemize}
    \end{block}
    \vspace{0.1cm}
    \begin{block}{Inputs \& Outputs}
      \textbf{Input:} Video frames, depth maps, GPS coordinates.\\
      \textbf{Output:} Directional spatial audio, speech synthesis.
    \end{block}
  \end{column}
  
  \begin{column}{0.31\textwidth}
    \begin{block}{Summary of Results}
      \begin{itemize}
        \item Vision-based systems outperform ultrasonic aids in semantic recognition.
        \item Edge-optimized models achieve real-time latency ($<50$\,ms).
      \end{itemize}
    \end{block}
    \vspace{0.1cm}
    \begin{alertblock}{Key Limitations}
      \begin{itemize}
        \item Most systems require specialized hardware rigs (wearables/vests).
        \item Lack of combined obstacle detection + semantic context.
      \end{itemize}
    \end{alertblock}
  \end{column}
\end{columns}
\end{frame}

% Paper 3: Manirajee et al. (SLR, 2024)
\begin{frame}{Literature Review: Systematic Review of Assistive Technologies}
\begin{columns}[T]
  \begin{column}{0.31\textwidth}
    \begin{block}{Paper Overview \cite{manirajee2024assistive}}
      \textbf{Authors:} L. Manirajee et al.\\
      \textbf{Journal:} IJARBSS (2024)\\
      \textbf{Focus:} Systematic Literature Review (SLR) on mobility and navigation aids for visually impaired.
    \end{block}
    \vspace{0.1cm}
    \begin{block}{Main Contributions}
      \begin{itemize}
        \item Systematic mapping of 50+ assistive technology publications.
        \item Analyzed cost vs. usability barrier across user demographics.
      \end{itemize}
    \end{block}
  \end{column}
    
  \begin{column}{0.34\textwidth}
    \begin{block}{Methodology \& Technology}
      \begin{itemize}
        \item Systematic PRISMA review methodology.
        \item Comparison of Electronic Travel Aids (ETAs), Electronic Orientation Aids (EOAs).
        \item AI computer vision and voice interaction frameworks.
      \end{itemize}
    \end{block}
    \vspace{0.1cm}
    \begin{block}{Inputs \& Outputs}
      \textbf{Input:} User mobility datasets, navigation metrics.\\
      \textbf{Output:} Comparative usability matrices.
    \end{block}
  \end{column}
  
  \begin{column}{0.31\textwidth}
    \begin{block}{Summary of Results}
      \begin{itemize}
        \item Identified smartphone solutions as the most scalable approach.
        \item Highlighted user preference for voice-guided navigation.
      \end{itemize}
    \end{block}
    \vspace{0.1cm}
    \begin{alertblock}{Key Limitations}
      \begin{itemize}
        \item Few existing works address developing nation road infrastructure.
        \item Absence of conversational AI for scene query answering.
      \end{itemize}
    \end{alertblock}
  \end{column}
\end{columns}
\end{frame}